%% SCRAP: archive/GAP_ANALYSIS
%% SOURCE: docs/working/archive/GAP_ANALYSIS.md
%% STATUS: HISTORICAL
%% FITS: none
%% EDITORIAL: lifted — prose rewritten to press voice

\section{Gap Analysis: Evolutionary Path to FPGA}

Conducted on 2025-12-14, this gap analysis documents the thirty-seven
critical components missing between the then-current StarForth VM and the
four-stage evolutionary target: StarForth hosted VM $\to$ StarKernel
(bare metal) $\to$ StarshipOS (self-hosting OS) $\to$ FPGA hardware
processor. The record is preserved as a historical planning artefact;
the majority of StarKernel gaps (Stages~1--2) have since been closed by
LithosAnanke~v1.0.8.

\subsection{State at Time of Analysis}

StarForth was a working FORTH-79 VM with a physics-driven adaptive runtime,
0\% algorithmic variance proven across ninety runs, and support for Linux
and L4Re. Documentation was strong: formal ontology, nine scientific
defence documents, HAL architecture documentation, and peer-review
materials. Build tooling included forty-plus Makefile targets, PGO
support, and a Design of Experiments framework.

\subsection{Gaps by Evolutionary Stage}

\subsubsection{Stage 1 — Hardware Abstraction Layer (Foundation)}

Eight critical gaps were identified (H1--H8): HAL headers for time,
interrupt, memory, console, CPU, and panic; a Linux HAL implementation;
and migration of the VM core to the HAL. The analysis concluded that
HAL implementation was the critical path: no progress toward StarKernel
or FPGA was possible without it. Estimated effort: two to three weeks.

\subsubsection{Stage 2 — StarKernel (Bare Metal Boot)}

Eight critical gaps (K1--K8) and five important gaps (K9--K13) covered
the UEFI boot loader, physical memory manager, virtual memory manager,
kernel heap allocator, GDT/IDT setup, UART driver, TSC/HPET time
sources, and APIC interrupt controller. Estimated source volume: eight
to ten thousand lines of C and assembly; estimated elapsed time to a
serial ``ok'' prompt: six to eight weeks.

\subsubsection{Stage 3 — StarshipOS (Full OS)}

Five critical gaps (OS1--OS5) addressed storage drivers, a filesystem,
block device abstraction, TCP/IP stack, and network drivers. Seven
important gaps (OS6--OS12) covered the process model, scheduler, IPC
mechanism, device model, security/capabilities, system call interface,
and userland tools. Estimated source volume: fifty thousand lines;
estimated elapsed time: twelve to eighteen months.

\subsubsection{Stage 4 — FPGA Implementation}

Eight critical gaps (F1--F8) and eleven additional gaps (F9--F19) spanned
feasibility study, HDL architecture, instruction set design, stack and
dictionary memory, memory controller, I/O IP cores, HDL language and
target FPGA selection, clock domain crossing, adaptive runtime in
hardware, synthesis constraints, and toolchain setup. The FPGA path
was described as entirely greenfield; estimated source volume: twenty to
thirty thousand lines of Verilog or VHDL.

\subsection{Critical Path}

The dependency graph placed HAL implementation (H1--H8) as the sole
gateway to all downstream work. Neither the StarKernel path nor the FPGA
path could begin until platform independence was established.

\subsection{Risk Assessment}

The highest technical-uncertainty risks were FPGA feasibility (F1) and
implementing the adaptive runtime in HDL (F12). The most resource-constrained
risk was the aggregate project scope: approximately ninety-three thousand
lines of new code representing twenty-eight to forty-two person-months of
solo effort.

\subsection{Total Gap Count}

\begin{itemize}
  \item Critical: 12
  \item Important: 15
  \item Nice-to-have: 10
  \item \textbf{Total: 37}
\end{itemize}

The analysis recommended that the FPGA path receive a formal go/no-go
feasibility decision (F1) before committing resources to HDL development.
